\section*{Introduction}
\label{sec:ce:intro}

Engineering a distributed cyber-physical system is an architecture decision before it is
a software decision. How many edge nodes and gateways to deploy, which hardware class
each should be, where each processing stage runs, how far the ingest stage is replicated
and which class of network link connects the tiers together determine whether the system
meets its deadline, what it costs, how much energy it draws and how much traffic it puts
on the network. These decisions interact: adding edge nodes shortens queues but raises
cost and idle power, moving analytics to the cloud unloads the edge but adds a wide-area
transfer to every request, and replication buys ingest throughput at the price of
compute load and synchronisation traffic.

The number of candidate architectures is the product of the decision domains and reaches
$10^{5}$--$10^{7}$ for even a modest set of variables. What makes this different from
classical processor-level design-space exploration is not the arithmetic but the price of
one answer: for distributed cyber-physical systems a single design point is evaluated by
simulation, and an industrial methodology built on execution traces of a lithography
machine reports up to a minute per design point while explicitly leaving the search
strategy out of scope\cite{saadatmand2025compdse}. A survey of the area concludes that
scalable exploration technology for such systems is "more or less
non-existing"\cite{herget2022dse}. The engineering consequence is concrete: the budget
that matters is the number of expensive evaluations an engineer can afford, not the size
of the space.

Two established options each give up something. Exhaustive enumeration returns the exact
set of optimal trade-offs but pays the evaluator for every candidate. Population-based
search such as NSGA-II\cite{deb2002nsga2} fits a fixed evaluation budget but returns an
approximation and, crucially for an engineering decision, cannot say which trade-offs it
missed. A third and smaller family reasons about \emph{partial} designs and prunes
regions that provably contain nothing optimal, either symbolically with constraint
solvers\cite{neubauer2018exact,lukasiewycz2008symbolic} or through bound-based Pareto
pruning in the wider optimisation
literature\cite{paretopruning2022,mandow2010namoa}. Integer programming achieves the same
guarantee for allocation problems in the edge--cloud
continuum\cite{kouloumpris2024optimization,kouloumpris2026reliability}.

Every method in that third family rests on one assumption: that each decision can only
add to the objective, so that a partially specified design already carries a valid lower
bound on all its completions. The dominant objective in a distributed cyber-physical
system violates it. End-to-end latency contains a waiting term that depends on the
utilisation of a tier, and utilisation is fixed only jointly, by the number of nodes,
their class, the replication factor and the placement policy at once. Adding a node
lowers it. There is therefore no monotone accumulation to bound, and the guarantee that
makes exact pruning attractive does not transfer.

Here we show that exact exploration is nevertheless achievable for these objectives, and
that it makes expensive architecture evaluation tractable rather than merely faster. We
construct admissible bounds for contention-dependent objectives by relaxing each additive
term of an objective independently over the domains still reachable, and by ordering the
decisions so that the determinants of tier utilisation are fixed early --- at which point
the waiting term stops being relaxed and becomes exactly computable on an incomplete
design. Pruning built on these bounds is provably safe, and the exploration returns the
exact Pareto front of the analytical model.

The reduction is large enough to change which studies are practical. On
architecture spaces of up to $8.3\times10^{7}$ candidates the method expands at most
$16\,780$ partial designs and calls the expensive evaluator $16$--$61$ times, while the
front it returns is identical, vector for vector, to the one exhaustive enumeration
returns wherever exhaustive enumeration is still feasible. An independent containerised
edge--gateway--cloud deployment, run with the analytical model frozen beforehand, shows
that the ordering the model produces is the ordering the hardware produces for
$88.6$--$91.3\%$ of the architecture pairs it separates by more than $25\%$ within a
workload. We report equally clearly where the approach does not help: when the evaluator
is cheap, when no admissible bound can be derived, and near saturation, where neither the
model nor the deployment is stable.
